\documentclass[11pt,a4paper]{article}

\usepackage[utf8]{inputenc}
\usepackage[T1]{fontenc}
\usepackage{amsmath,amssymb}
\usepackage{graphicx}
\usepackage{booktabs}
\usepackage{hyperref}
\usepackage[margin=1in]{geometry}
\usepackage{listings}
\usepackage{xcolor}
\usepackage{fancyhdr}
\usepackage{titlesec}
\usepackage{enumitem}

% Code listing style
\lstset{
    language=Python,
    basicstyle=\ttfamily\small,
    keywordstyle=\color{blue}\bfseries,
    stringstyle=\color{red!70!black},
    commentstyle=\color{green!50!black}\itshape,
    numbers=left,
    numberstyle=\tiny\color{gray},
    numbersep=8pt,
    frame=single,
    framesep=4pt,
    breaklines=true,
    breakatwhitespace=true,
    showstringspaces=false,
    tabsize=4,
    backgroundcolor=\color{gray!5},
    rulecolor=\color{gray!30},
}

\pagestyle{fancy}
\fancyhead[L]{\small BioMind-BG Technical Documentation}
\fancyhead[R]{\small \thepage}
\fancyfoot{}

\title{\textbf{BioMind-BG: Complete Technical Documentation}\\[0.5em]
\large A Biologically Faithful Basal Ganglia Simulation\\
File-by-File Code Reference with Biological Explanations}

\author{BioMind Project}
\date{March 2026}

\begin{document}

\maketitle
\tableofcontents
\newpage

% ============================================================================
% OVERVIEW
% ============================================================================
\section{Project Overview}

BioMind-BG is a spiking neural network simulation of the cortico-basal-ganglia-thalamic (CBGT) circuit. It simulates approximately 3,000 neurons organized into 9 brain nuclei, connected by 27 biologically constrained pathways, learning from dopamine reward signals.

\subsection{What It Does}
\begin{enumerate}
    \item Builds 9 populations of neurons (Cortex, D1-MSN, D2-MSN, FSI, GPe, STN, GPi, Thalamus, Cortical Interneurons)
    \item Wires them together exactly as they are in a real brain
    \item Simulates spiking dynamics with AMPA, GABA, and NMDA receptors
    \item Detects decisions via thalamic firing rate threshold crossing
    \item Learns from reward through dopamine-modulated synaptic plasticity
\end{enumerate}

\subsection{File Structure}
\begin{table}[h]
\centering
\begin{tabular}{lll}
\toprule
File & Lines & Purpose \\
\midrule
\texttt{params.py} & 190 & All biological constants with units \\
\texttt{populations.py} & 130 & Population construction and connectivity matrices \\
\texttt{agent.py} & 160 & Neural state initialization (all arrays) \\
\texttt{timestep.py} & 260 & Differential equations (the core simulation loop) \\
\texttt{qlearning.py} & 65 & Q-values and RPE-to-dopamine conversion \\
\texttt{trial.py} & 130 & Trial state machine (stimulus, decision, reward) \\
\texttt{run.py} & 170 & Main simulation runner \\
\bottomrule
\end{tabular}
\end{table}

\subsection{Data Flow}
The data flows through the system in this order:
\begin{enumerate}
    \item \texttt{params.py} defines constants
    \item \texttt{populations.py} builds population dicts and connectivity matrices from those constants
    \item \texttt{agent.py} expands everything into per-neuron NumPy arrays
    \item \texttt{timestep.py} runs the differential equations every 0.2ms
    \item \texttt{trial.py} manages the behavioral task (stimulus, decision detection, reward delivery)
    \item \texttt{qlearning.py} converts reward outcomes into dopamine signals
    \item \texttt{run.py} orchestrates everything
\end{enumerate}

\newpage
% ============================================================================
% PARAMS.PY
% ============================================================================
\section{params.py --- All Biological Constants}

This file contains every number the simulation uses. No computation happens here, just definitions. Every parameter has a unit and a biological justification.

\subsection{Simulation Constants}

\begin{lstlisting}
DT = 0.2          # ms, integration timestep
BUFFER_LEN = 300  # rolling window for firing rate
\end{lstlisting}

\textbf{DT = 0.2 ms:} A real spike lasts about 1--2 ms. If our timestep were larger than that, we would miss spikes entirely. 0.2 ms gives us 5--10 samples per spike, which is enough to track the voltage dynamics accurately. Smaller would be more precise but slower. This is the sweet spot that CBGTPy validated.

\textbf{BUFFER\_LEN = 300:} The rolling buffer for measuring firing rates. 300 slots $\times$ 0.2 ms = 60 ms window. A neuron fires spikes as discrete events. To measure ``firing rate'' we need to count spikes over some window of time. 60 ms is long enough for a stable average but short enough to track rapid changes during a decision. The buffer is circular: when slot 300 is reached, it wraps back to slot 0.

\subsection{Default Neuron Parameters}

Every neuron in the model starts with these defaults. Specific populations then override what they need.

\begin{lstlisting}
NEURON_DEFAULTS = {
    'N':          75,      # neurons per population per channel
    'C':          0.5,     # nF, membrane capacitance
    'Taum':       20.0,    # ms, membrane time constant
    'RestPot':    -70.0,   # mV, resting membrane potential
    'ResetPot':   -55.0,   # mV, post-spike reset potential
    'Threshold':  -50.0,   # mV, spike threshold
}
\end{lstlisting}

\textbf{N = 75:} Each population gets 75 neurons per action channel. Not 1, not 10,000. 75 is enough to get population-level statistics (averaging over noise) but small enough to simulate fast.

\textbf{C = 0.5 nF:} Capacitance. The cell membrane is a thin insulator between two conductors (inside and outside the cell). It stores charge like a capacitor. Higher C means the neuron needs more current to change its voltage. Think of it as the neuron's ``inertia.''

\textbf{Taum = 20 ms:} Membrane time constant. This is how fast voltage decays back to rest when there is no input. At 20 ms, the neuron ``forgets'' an input after about 20 ms. This sets the timescale of integration.

\textbf{RestPot = $-$70 mV:} Where the voltage sits when nothing is happening. This is set by the balance of leak K$^+$ and Na$^+$ channels. Every neuron in your brain right now is sitting near $-$70 mV.

\textbf{ResetPot = $-$55 mV and Threshold = $-$50 mV:} After a spike, voltage resets to $-$55. The next spike requires reaching $-$50. That is only a 5 mV gap, which means neurons can fire again quickly if input is strong enough.

\subsubsection{Calcium Dynamics}

\begin{lstlisting}
    'RestPot_ca': -85.0,   # mV, calcium reversal potential
    'Alpha_ca':   0.5,     # uM, calcium increment per spike
    'Tau_ca':     80.0,    # ms, calcium decay time constant
    'Eff_ca':     0.0,     # calcium efficacy (unused)
\end{lstlisting}

Every time a neuron spikes, a little calcium flows in (Alpha\_ca = 0.5). This calcium activates a potassium channel (AHP, afterhyperpolarization) that pulls the voltage \textbf{down}. So after a burst of spikes, calcium builds up and the neuron gets harder to fire. It is a self-braking mechanism --- spike-rate adaptation. The calcium then slowly decays away with Tau\_ca = 80 ms, letting the neuron recover.

\subsubsection{T-Current (Rebound Bursting)}

\begin{lstlisting}
    'tauhm':      20.0,    # ms, burst inactivation time constant
    'tauhp':      100.0,   # ms, de-inactivation time constant
    'V_h':        -60.0,   # mV, burst activation threshold
    'V_T':        120.0,   # mV, T-current reversal potential
    'g_T':        0.0,     # mS/cm2, max conductance (OFF by default)
\end{lstlisting}

This mechanism is \textbf{OFF by default} ($g_T = 0$) but turned ON for STN and GPe ($g_T = 0.06$). When a neuron is held below $-$60 mV (inhibited) for a while, the $h$ gating variable slowly charges up (de-inactivation, $\tau = 100$ ms). When the inhibition releases, that stored $h$ fires a burst of calcium current that pushes the neuron way above threshold. That is a rebound burst.

This is how STN fires a burst when GPe releases it from inhibition --- critical for the indirect pathway.

\subsubsection{Additional Channels}

\begin{lstlisting}
    'g_adr_max':  0.0,     # Anomalous delayed rectifier (OFF)
    'g_k_max':    0.0,     # Outward rectifying K+ (OFF)
\end{lstlisting}

Both are OFF by default. Included for biological completeness and future extensions.

\subsection{Population-Specific Overrides}

\begin{lstlisting}
POP_SPECIFIC = {
    'Cx':   {'N': 204, 'dpmn_cortex': 1},
    'CxI':  {'N': 186, 'C': 0.2, 'Taum': 10.0},
    'dSPN': {},
    'iSPN': {},
    'FSI':  {'C': 0.2, 'Taum': 10.0},
    'GPe':  {'N': 750, 'g_T': 0.06, 'Taum': 20.0},
    'STN':  {'N': 750, 'g_T': 0.06},
    'GPi':  {},
    'Th':   {'Taum': 27.78},
}
\end{lstlisting}

\begin{table}[h]
\centering
\caption{Population overrides and biological justification.}
\begin{tabular}{lll}
\toprule
Population & Overrides & Why \\
\midrule
Cx (Cortex) & N=204, dpmn\_cortex=1 & Largest pop. Flag for plasticity source \\
CxI (Cortical Interneurons) & C=0.2, Taum=10 & PV+ cells: small, fast \\
dSPN (D1-MSN) & \{\} (all defaults) & Quiet neurons, need strong input \\
iSPN (D2-MSN) & \{\} (all defaults) & Same as dSPN \\
FSI (Fast-Spiking) & C=0.2, Taum=10 & PV+ interneurons: small, fast \\
GPe & N=750, g\_T=0.06 & Pacemaker, rebound bursting ON \\
STN & N=750, g\_T=0.06 & Rebound bursting ON \\
GPi & \{\} (all defaults) & Output gate \\
Th (Thalamus) & Taum=27.78 & Slower integration, smooths GPi input \\
\bottomrule
\end{tabular}
\end{table}

\subsection{Receptor Kinetics}

\begin{lstlisting}
RECEPTOR_DEFAULTS = {
    'Tau_AMPA':     2.0,    # ms
    'RevPot_AMPA':  0.0,    # mV
    'Tau_GABA':     5.0,    # ms
    'RevPot_GABA':  -70.0,  # mV
    'Tau_NMDA':     100.0,  # ms
    'RevPot_NMDA':  0.0,    # mV
}
NMDA_ALPHA = 0.6332
\end{lstlisting}

\textbf{AMPA} ($\tau = 2$ ms, $E_{rev} = 0$ mV): Fast excitatory. Opens immediately on glutamate binding, lets Na$^+$ in, pushes voltage toward 0 mV. The only plastic receptor --- all learning happens at AMPA synapses.

\textbf{GABA$_A$} ($\tau = 5$ ms, $E_{rev} = -70$ mV): Fast inhibitory. Opens Cl$^-$ channels, clamps voltage at $-$70 mV (rest). This is how dSPN inhibits GPi, how GPi inhibits Thalamus, how GPe inhibits STN.

\textbf{NMDA} ($\tau = 100$ ms, $E_{rev} = 0$ mV): Slow excitatory. Same reversal as AMPA but 50$\times$ slower. A single spike's effect lingers for 100 ms. Has two special properties: Mg$^{2+}$ block (only opens when neuron is already depolarized) and saturation (NMDA\_ALPHA prevents unbounded accumulation).

\subsection{External Background Input}

\begin{lstlisting}
BASESTIM = {
    'Cx':   {'FreqExt_AMPA': 2.3,  'MeanExtEff_AMPA': 2.0,
             'MeanExtCon_AMPA': 800},
    'GPe':  {'FreqExt_AMPA': 4.0,  'MeanExtEff_AMPA': 2.0,
             'MeanExtCon_AMPA': 800,
             'FreqExt_GABA': 2.0,  'MeanExtEff_GABA': 2.0,
             'MeanExtCon_GABA': 2000},
    ...
}
\end{lstlisting}

Every nucleus is constantly bombarded by thousands of synapses from brain regions we are not modeling. Three numbers per population work together: \textbf{FreqExt} (how fast external neurons fire), \textbf{MeanExtEff} (how strong each connection is), \textbf{MeanExtCon} (how many external connections). The total background drive is roughly Freq $\times$ Eff $\times$ N.

Key observation: \textbf{GPe is the only population with external GABA input} (2000 connections). This represents striatal inhibition from brain regions outside our model.

\subsection{Pathway Connectivity}

The complete wiring diagram. Each row: (source, destination, receptor, type, connection probability, efficacy, plastic).

\textbf{Type} is critical: \texttt{'syn'} means channel-specific (action 1 cortex only connects to action 1 striatum). \texttt{'all'} means all-to-all across channels (STN excites ALL GPi channels).

\begin{table}[h]
\centering
\caption{Key pathways in the CBGT circuit.}
\small
\begin{tabular}{llllccc}
\toprule
Source & Dest & Receptor & Type & Prob & Eff & Plastic \\
\midrule
\multicolumn{7}{l}{\textbf{Direct Pathway (Go)}} \\
Cx & dSPN & AMPA & syn & 1.0 & 0.015 & \textbf{Yes} \\
dSPN & GPi & GABA & syn & 1.0 & 2.09 & No \\
\midrule
\multicolumn{7}{l}{\textbf{Indirect Pathway (NoGo)}} \\
Cx & iSPN & AMPA & syn & 1.0 & 0.015 & \textbf{Yes} \\
iSPN & GPe & GABA & syn & 1.0 & 4.07 & No \\
GPe & STN & GABA & syn & 0.067 & 0.35 & No \\
STN & GPi & NMDA & all & 1.0 & 0.038 & No \\
\midrule
\multicolumn{7}{l}{\textbf{Output}} \\
GPi & Th & GABA & syn & 1.0 & 0.3315 & No \\
\bottomrule
\end{tabular}
\end{table}

Only \textbf{two} connections are plastic: Cx$\rightarrow$dSPN AMPA and Cx$\rightarrow$iSPN AMPA. All learning in this model happens at the corticostriatal synapse, modulated by dopamine.

\subsection{Dopamine Parameters}

\begin{lstlisting}
D1_PARAMS = {
    'dpmn_type':    1,
    'dpmn_alphaw':  39.5,     # positive = LTP on reward
    'dpmn_wmax':    0.055,
}
D2_PARAMS = {
    'dpmn_type':    2,
    'dpmn_alphaw':  -38.2,    # negative = LTD on reward
    'dpmn_wmax':    0.035,
}
\end{lstlisting}

The sign of \texttt{alphaw} is the entire Go/NoGo mechanism. When dopamine fires (reward): D1 gets $+39.5 \times$ positive fDA = LTP (Go strengthened). D2 gets $-38.2 \times$ positive fDA = LTD (NoGo weakened). On punishment, everything reverses.

\subsection{Q-Learning Parameters}

\begin{lstlisting}
Q_PARAMS = {
    'q_alpha':  0.1,    # 10% update per trial
    'C_scale':  80.0,   # RPE amplification for neural scale
    'Q_init':   0.5,    # no prior bias
}
\end{lstlisting}

\texttt{C\_scale = 80} converts an abstract prediction error (small number between $-$2 and $+$1) into a neural-scale dopamine signal. Without this scaling, the RPE would produce negligible weight changes in the spiking network.

\subsection{Trial Parameters}

\begin{lstlisting}
TRIAL_DEFAULTS = {
    'thalamic_threshold':  30.0,    # Hz
    'choice_timeout':      1000,    # ms
    'movement_time':       300,     # ms
    'inter_trial_interval': 300,    # ms
    'max_stim':            3.0,
    'warmup_steps':        5000,
}
\end{lstlisting}

\textbf{thalamic\_threshold = 30 Hz:} A decision happens when any thalamic channel exceeds 30 Hz. At baseline, thalamus fires at $\sim$7 Hz. It must roughly quadruple. This only happens when dSPN successfully inhibits GPi enough to release thalamus.

\textbf{movement\_time = 300 ms:} Post-decision delay before reward. The eligibility trace ($\tau_E = 100$ ms) must survive this delay for learning to work.

\textbf{max\_stim = 3.0:} During stimulus, cortex gets baseline + 3.0 kHz, driving it from $\sim$2.3 to 5.3 kHz, producing $\sim$116 Hz firing.

\newpage
% ============================================================================
% POPULATIONS.PY
% ============================================================================
\section{populations.py --- Building the Network}

This file answers two questions: (1) How do we create actual populations from the parameter templates? (2) How do we build the weight matrices connecting them?

\subsection{Channel Architecture}

For a 2-choice task, some populations are \textbf{channel-specific} (each action gets its own copy) and some are \textbf{shared} (one copy serves all actions).

\begin{itemize}
    \item \textbf{Channel-specific:} GPi, STN, GPe, dSPN, iSPN, Cx, Th. Each action needs its own set. Action 0's cortex does not leak into action 1's striatum.
    \item \textbf{Shared:} FSI, CxI. Interneurons that provide blanket inhibition across all channels.
\end{itemize}

For 2 actions: $7 \times 2 + 2 = 16$ populations total.

\subsection{build\_population\_data()}

\begin{lstlisting}
def build_population_data(n_actions=2):
    channel_pops = ['GPi','STN','GPe','dSPN','iSPN','Cx','Th']
    shared_pops = ['FSI', 'CxI']
    pops = []
    pop_index = {}
    for name in POP_NAMES:
        if name in channel_pops:
            for ch in range(n_actions):
                pop = _build_single_pop(name, ch)
                idx = len(pops)
                pops.append(pop)
                pop_index[(name, ch)] = idx
        else:
            pop = _build_single_pop(name, -1)
            idx = len(pops)
            pops.append(pop)
            pop_index[(name, -1)] = idx
    return pops, pop_index
\end{lstlisting}

\texttt{pop\_index} is a lookup dictionary. To find ``where is dSPN for action 0?'' you ask \texttt{pop\_index[('dSPN', 0)]} and get back the array index. Shared populations use channel = $-$1.

\subsection{\_build\_single\_pop()}

Assembles a complete parameter dictionary by layering:
\begin{enumerate}
    \item Start with NEURON\_DEFAULTS (the template)
    \item Apply POP\_SPECIFIC overrides (e.g., GPe gets $g_T = 0.06$)
    \item Add RECEPTOR\_DEFAULTS (AMPA/GABA/NMDA kinetics)
    \item Add BASESTIM (background bombardment)
    \item Fill missing stim params with 0 (prevents KeyError)
    \item Add DPMN params only for dSPN and iSPN
\end{enumerate}

\subsection{build\_connectivity()}

Turns the PATHWAYS table into actual NumPy weight matrices.

\begin{lstlisting}
def build_connectivity(pops, pop_index, n_actions=2):
    for receptor in ['AMPA', 'GABA', 'NMDA']:
        for pathway in PATHWAYS:
            if rec != receptor: continue
            pairs = _get_population_pairs(...)
            for src_idx, dest_idx in pairs:
                if prob >= 1.0:
                    con_data = np.ones((n_src, n_dest))
                else:
                    con_data = (np.random.rand(n_src,n_dest) < prob)
                eff_data = con_data * eff
\end{lstlisting}

For each population pair: if connection probability is 1.0, all synapses exist. If less than 1.0, each possible synapse gets a random draw. \textbf{Efficacy = connection $\times$ weight}. If a synapse does not exist (con = 0), its efficacy is automatically 0.

\subsection{\_get\_population\_pairs() --- The syn vs all Logic}

\texttt{'syn'}: Only matching channels connect. Cx channel 0 $\rightarrow$ dSPN channel 0. No cross-talk. The condition is \texttt{src\_ch == dest\_ch} (or either is $-$1 for shared populations).

\texttt{'all'}: Every source connects to every destination regardless of channel. This is how STN provides a broad brake signal to all GPi channels simultaneously.

\newpage
% ============================================================================
% AGENT.PY
% ============================================================================
\section{agent.py --- The Living Network State}

\texttt{params.py} is the blueprint. \texttt{populations.py} is the construction plan. \texttt{agent.py} is the actual building with all its rooms and plumbing ready to run.

Every attribute is a list of NumPy arrays indexed by population ID. So \texttt{self.V[3]} is the voltage array for population 3, and \texttt{self.V[3][47]} is the voltage of neuron 47 in that population.

\subsection{Core State Variables}

\begin{lstlisting}
self.V = [np.ones(p['N']) * p['RestPot'] for p in pops]
self.Ca = [np.zeros(p['N']) for p in pops]
self.h = [np.ones(p['N']) * p['h'] for p in pops]
self.RefrState = [np.zeros(p['N'], dtype=int) for p in pops]
self.spikes = [[] for _ in pops]
self.timesincelastspike = [np.zeros(p['N']) for p in pops]
self.Ptimesincelastspike = [np.zeros(p['N']) for p in pops]
\end{lstlisting}

\begin{itemize}
    \item \textbf{V} starts at $-$70 mV for every neuron. Resting state.
    \item \textbf{Ca} starts at 0. Calcium only accumulates when neurons spike.
    \item \textbf{h} starts at 1.0 (fully de-inactivated). STN and GPe are ready to rebound burst immediately.
    \item \textbf{RefrState} starts at 0. When a neuron spikes, set to 10 (= 2 ms refractory).
    \item \textbf{spikes} is rebuilt every timestep. Contains indices of neurons that fired.
    \item \textbf{Ptimesincelastspike} saves the timer value BEFORE reset, needed for NMDA saturation.
\end{itemize}

\subsection{Parameter Expansion}

Every scalar parameter gets expanded to a per-neuron array using \texttt{np.full()}:
\begin{lstlisting}
self.Taum = [np.full(p['N'], p['Taum']) for p in pops]
\end{lstlisting}

This creates an array of 75 identical values. Why not use the scalar directly? Because the timestep code does \texttt{a.V[i] += ... / a.Taum[i]} where \texttt{V[i]} is a 75-element array. Keeping everything as arrays is consistent and enables future per-neuron heterogeneity.

\subsection{Synaptic Conductance State}

Two sources of synaptic input per receptor:

\begin{itemize}
    \item \textbf{ExtS} (external): Background noise following Ornstein-Uhlenbeck process. Fluctuates around \textbf{ExtMuS} with standard deviation \textbf{ExtSigmaS}.
    \item \textbf{LS} (lateral): Driven by spikes from other populations in the network. When neuron $j$ in GPe spikes, it adds to \texttt{LS\_GABA} of connected STN neurons.
\end{itemize}

\subsection{Connectivity Matrices}

Unpacked from the connectivity dict built by \texttt{populations.py}:
\begin{lstlisting}
self.AMPA_con, self.AMPA_eff, self.AMPA_plastic = connectivity['AMPA']
self.GABA_con, self.GABA_eff, _ = connectivity['GABA']
self.NMDA_con, self.NMDA_eff, _ = connectivity['NMDA']
\end{lstlisting}

Only AMPA plasticity is kept. GABA and NMDA plasticity matrices are discarded because no GABA or NMDA synapses are plastic.

\textbf{LastConductanceNMDA} tracks NMDA saturation: a 2D structure mirroring the connection matrices, recording the conductance level at each synapse's last spike.

\subsection{Dopamine State}

All populations get dopamine arrays, but only dSPN (\texttt{dpmn\_type} = 1) and iSPN (\texttt{dpmn\_type} = 2) are actually processed. For all other populations, \texttt{dpmn\_type} = 0 and the plasticity code skips them.

Key dynamic variables:
\begin{itemize}
    \item \textbf{dpmn\_DAp}: Phasic dopamine. Set by reward, decays with $\tau = 2$ ms.
    \item \textbf{dpmn\_APRE/APOST}: Running spike traces for STDP.
    \item \textbf{dpmn\_E}: Eligibility trace. Bridges STDP timing to dopamine arrival.
    \item \textbf{dpmn\_XPRE/XPOST}: Binary spike indicators. Reset each timestep.
\end{itemize}

\subsection{Temporary Arrays and Firing Rate Buffer}

Pre-allocated scratch pads (\texttt{cond}, \texttt{g\_rb}, \texttt{g\_adr}, \texttt{g\_k}, \texttt{dv}, \texttt{tau\_n}, \texttt{n\_inif}, \texttt{Vaux}) avoid memory allocation every timestep. \texttt{cond} is the integer mask: 1 for neurons to update, 0 for neurons in refractory or above threshold.

The \textbf{rollingbuffer} is the circular firing rate buffer: \texttt{n\_pops $\times$ 300} array. Each timestep writes spike count into the current column, pointer advances and wraps.

\textbf{get\_firing\_rates()} computes: mean spikes per timestep $\div$ N neurons $\div$ dt $\times$ 1000 = Hz.

\newpage
% ============================================================================
% TIMESTEP.PY
% ============================================================================
\section{timestep.py --- The Core Simulation}

This is the heart of BioMind-BG. Every other file sets up data structures. This file makes neurons \textbf{fire}. It computes one $dt = 0.2$ ms step. Called 5000 times per second of brain time.

\subsection{Step 1: Reset Spike Indicators}

\begin{lstlisting}
for i in range(n_pops):
    a.dpmn_XPRE[i] *= 0
    a.dpmn_XPOST[i] *= 0
\end{lstlisting}

Every timestep starts fresh. XPRE and XPOST are binary flags set during Steps 3 and 6 when spikes occur.

\subsection{Step 2: External Input (Ornstein-Uhlenbeck)}

\begin{lstlisting}
a.ExtMuS_AMPA[i] = eff * freq * 0.001 * N_conn * tau
a.ExtSigmaS_AMPA[i] = eff * sqrt(tau * 0.5 * freq * 0.001 * N_conn)
a.ExtS_AMPA[i] += dt/tau * (-ExtS + Mu) 
                 + Sigma * sqrt(2*dt/tau) * N(0,1)
a.LS_AMPA[i] *= exp(-dt / tau)
\end{lstlisting}

The external conductance follows a random walk pulled back toward its mean (Ornstein-Uhlenbeck process). Like a drunk person on a leash:
\begin{itemize}
    \item First term: \texttt{dt/tau * (-ExtS + Mu)} pulls conductance toward the mean.
    \item Second term: Gaussian noise provides random kicks.
    \item Lateral conductance decays exponentially, waiting for new spikes in Step 3.
\end{itemize}

\subsection{Step 3: Spike Propagation}

For every neuron that spiked last timestep, add its synaptic weight to destination conductances:

\textbf{AMPA:} Direct increment + XPRE tracking for plasticity.

\textbf{GABA:} Direct increment only.

\textbf{NMDA:} Includes saturation mechanism. Each spike adds $\alpha \times (1 - \text{last\_conductance})$ instead of a fixed amount, preventing unbounded accumulation during sustained firing.

This is the \textbf{slowest part of the code} due to the triple nested Python loop (populations $\times$ populations $\times$ spiking neurons).

\subsection{Step 4: T-Current (Rebound Bursting)}

\begin{lstlisting}
cond = (a.V[i] < a.V_h[i]).astype(int)
a.h[i] += cond * dt * (1.0 - a.h[i]) / a.tauhp[i]       # charge
a.h[i] += (1 - cond) * dt * (-a.h[i]) / a.tauhm[i]       # decay
a.g_rb[i] = a.g_T[i] * a.h[i] * (1 - cond)               # burst
\end{lstlisting}

When hyperpolarized ($V < -60$ mV): $h$ charges toward 1.0 (100 ms). When released: $h$ drives a burst through $g_T$, then decays (20 ms). Only active in STN and GPe ($g_T = 0.06$). All other populations: $g_T = 0$, so $g_{rb} = 0$ always.

\subsection{Step 5: Membrane Voltage Update}

The main equation with five current terms:

\begin{equation}
\frac{dV}{dt} = -\left[\underbrace{\frac{V - V_{rest}}{\tau_m}}_{\text{leak}} + \underbrace{\frac{Ca \cdot g_{AHP}}{C}(V - V_K)}_{\text{afterhyperpol.}} + \underbrace{\frac{g_{ADR}}{C}(V - V_{ADR})}_{\text{delayed rect.}} + \underbrace{\frac{g_K}{C}(V - V_{ADR})}_{\text{outward K}^+} + \underbrace{\frac{g_{rb}}{C}(V - V_T)}_{\text{T-current}}\right]
\end{equation}

Then three synaptic currents are added:

\begin{align}
\Delta V_{\text{AMPA}} &= dt \cdot \frac{(E_{\text{AMPA}} - V) \cdot (S^{lat}_{\text{AMPA}} + S^{ext}_{\text{AMPA}})}{C} \times 0.001 \\
\Delta V_{\text{GABA}} &= dt \cdot \frac{(E_{\text{GABA}} - V) \cdot (S^{lat}_{\text{GABA}} + S^{ext}_{\text{GABA}})}{C} \times 0.001 \\
\Delta V_{\text{NMDA}} &= dt \cdot \frac{(E_{\text{NMDA}} - V) \cdot (S^{lat}_{\text{NMDA}} + S^{ext}_{\text{NMDA}})}{C \cdot (1 + e^{-0.062 V / 3.57})} \times 0.001
\end{align}

The $(E_{rev} - V)$ term is the driving force. For AMPA ($E = 0$ mV): at rest ($V = -70$), driving force = $+70$ (excitatory). For GABA ($E = -70$ mV): at threshold ($V = -50$), driving force = $-20$ (inhibitory). NMDA has the additional Mg$^{2+}$ block denominator.

Neurons in refractory period (\texttt{cond = 0}) skip the entire update.

\subsection{Step 6: Spike Detection}

\begin{lstlisting}
spiked = np.nonzero(a.V[i] > a.Threshold[i])[0]
for neuron in spiked:
    a.V[i][neuron] = 0.0           # spike marker
    a.Ca[i][neuron] += alpha_ca    # calcium influx
    a.RefrState[i][neuron] = 10    # 2ms refractory
    a.Ptimesincelastspike = a.timesincelastspike  # save
    a.timesincelastspike = 0.0     # reset timer
    a.dpmn_XPOST[i][neuron] = 1.0  # mark for plasticity
\end{lstlisting}

Six things happen per spike: voltage marked, calcium added, refractory started, timer saved then reset, post-synaptic indicator set.

\subsection{Step 7: Dopamine Plasticity (Three-Factor Learning)}

Only runs for populations with \texttt{dpmn\_type > 0} (dSPN and iSPN).

\textbf{Factor 1 --- STDP Eligibility Trace:}
\begin{align}
\frac{dA_{pre}}{dt} &= \frac{\delta_{pre} \cdot X_{pre} - A_{pre}}{\tau_{pre}} \\
\frac{dA_{post}}{dt} &= \frac{\delta_{post} \cdot X_{post} - A_{post}}{\tau_{post}} \\
\frac{dE}{dt} &= \frac{X_{post} \cdot A_{pre} - X_{pre} \cdot A_{post} - E}{\tau_E}
\end{align}

$E > 0$: pre-before-post (LTP direction). $E < 0$: post-before-pre (LTD direction). $\tau_E = 100$ ms holds this signal for dopamine arrival.

\textbf{Factor 2 --- Dopamine Signal:}
\begin{equation}
DA = m \cdot (DA_{phasic} + DA_{tonic})
\end{equation}

Transformed through receptor-specific $f(DA)$:
\begin{itemize}
    \item D1: linear with slope $y/x = 6.0$, clamps at $-y = -3.0$ below $DA = -0.5$
    \item D2: linear with slope $y/x \cdot \varepsilon = 1.8$, clamps at $y \cdot \varepsilon = 0.9$ above $DA = 0.5$
\end{itemize}

\textbf{Factor 3 --- Weight Update:}
\begin{equation}
\Delta w = dt \cdot c_{ij} \cdot \alpha_w \cdot f(DA) \cdot E \cdot \begin{cases} (w_{max} - w) & \text{if } \Delta w > 0 \\ (w - 0.001) & \text{if } \Delta w < 0 \end{cases}
\end{equation}

Soft multiplicative bounds: weights approach but never reach $w_{max}$ or 0.001.

\subsection{Step 8: Firing Rate Buffer}

\begin{lstlisting}
for i in range(n_pops):
    a.rollingbuffer[i][a.bufferpointer] = len(a.spikes[i])
a.bufferpointer = (a.bufferpointer + 1) % a.bufferlength
\end{lstlisting}

Write spike count, advance pointer, wrap at 300.

\newpage
% ============================================================================
% QLEARNING.PY
% ============================================================================
\section{qlearning.py --- Reward to Dopamine Bridge}

This module bridges \textbf{behavior} and \textbf{neurons}. It translates ``you chose action 0 and got $+$1 reward'' into a dopamine signal that the spiking network can use to update its weights.

It implements the Rescorla-Wagner learning rule (1972), which is the simplest form of reinforcement learning. The prediction error maps to Schultz's (1997) RPE theory of dopamine: dopamine neurons fire when outcomes are \textbf{better} than expected, pause when \textbf{worse} than expected.

\subsection{QLearner Class}

\begin{lstlisting}
class QLearner:
    def __init__(self, n_actions=2):
        self.alpha = 0.1       # learning rate
        self.C_scale = 80.0    # RPE amplification
        self.Q = np.full(n_actions, 0.5)  # initial Q-values
        self.Q_history = [self.Q.copy()]
        self.DA_history = []
\end{lstlisting}

\textbf{alpha = 0.1}: Q-values update by 10\% of the prediction error per trial. Slow, stable learning.

\textbf{C\_scale = 80}: Amplifies the abstract RPE (a number between $-$2 and $+$1) to neural scale. RPE of $+$0.5 becomes DA of $+$40.

\textbf{Q starts at 0.5}: No prior bias. Both actions equally valued.

\subsection{update() --- The Core Learning Step}

Three things happen every trial:

\begin{enumerate}
    \item \textbf{Prediction Error:} $RPE = reward - Q[chosen]$
    
    Example: $Q[0] = 0.5$, reward $= +1.0 \Rightarrow RPE = +0.5$ (better than expected)
    
    \item \textbf{Q-Value Update:} $Q[chosen] \mathrel{+}= \alpha \cdot RPE$
    
    Example: $Q[0] = 0.5 + 0.1 \times 0.5 = 0.55$
    
    \item \textbf{Dopamine Signal:} $DA = RPE \times C_{scale}$
    
    Example: $DA = 0.5 \times 80 = +40.0$ (injected into all striatal neurons)
\end{enumerate}

\newpage
% ============================================================================
% TRIAL.PY
% ============================================================================
\section{trial.py --- The Behavioral Task}

Manages the three-phase trial loop that sits on top of the neural simulation.

\subsection{Phase 0: Stimulus Presentation}

Cortical input is turned on. Each action channel gets \texttt{FreqExt\_AMPA = basestim + gain $\times$ max\_stim}. The thalamic firing rate is monitored every millisecond (5 timesteps).

A \textbf{decision} is detected when any thalamic channel crosses 30 Hz. The first channel to cross wins (first-past-the-post). If 1000 ms pass with no crossing, it is a timeout (no choice).

The neural mechanism: stimulus drives cortex $\rightarrow$ cortex excites dSPN $\rightarrow$ dSPN inhibits GPi $\rightarrow$ GPi's tonic inhibition of thalamus weakens $\rightarrow$ thalamus firing rate rises $\rightarrow$ crosses 30 Hz $\rightarrow$ decision.

\subsection{Phase 1: Movement Delay}

After decision, gain on the unchosen action drops to 0, chosen action drops to 30\%. Network waits 300 ms. During this delay, the eligibility trace ($\tau_E = 100$ ms) is decaying. Dopamine must arrive before it fades.

\subsection{Phase 2: Reward and Inter-Trial Interval}

The chosen action is evaluated against the reward schedule. Q-values update. Phasic dopamine is computed and injected into ALL striatal populations (both dSPN and iSPN, both channels). After 300 ms ITI, everything resets for the next trial.

\subsection{Key Implementation Detail}

\begin{lstlisting}
def step(self):
    # Apply stimulus to cortex
    # Run 5 neural timesteps (1 ms)
    from biomind.timestep import multi_timestep
    multi_timestep(a, 5)
    # Check for phase transitions
\end{lstlisting}

Each call to \texttt{step()} advances the simulation by 1 ms (5 $\times$ 0.2 ms). The trial manager calls this in a loop until all trials are complete.

\newpage
% ============================================================================
% RUN.PY
% ============================================================================
\section{run.py --- Orchestration}

Ties everything together in a simple sequence:

\begin{enumerate}
    \item Build populations and connectivity (\texttt{populations.py})
    \item Initialize agent (\texttt{agent.py})
    \item Zero cortical stimulus and run 5000-step warmup (\texttt{timestep.py})
    \item Generate reward schedule (probabilistic, from seed)
    \item Create QLearner and TrialManager
    \item Loop: call \texttt{trial\_mgr.step()} until all trials complete
    \item Collect and return results
\end{enumerate}

\texttt{run\_baseline\_test()} runs only the warmup and reports firing rates, useful for validating the network without running trials.

\texttt{run\_simulation()} runs the full experiment with trials, learning, and behavioral output.

\newpage
% ============================================================================
% VALIDATED RESULTS
% ============================================================================
\section{Validated Results}

\subsection{Experiment 1: Baseline Firing Rates}

\begin{table}[h]
\centering
\begin{tabular}{lccc}
\toprule
Population & Model (Hz) & Expected (Hz) & Status \\
\midrule
Cortex & 0.0 & 0--5 & PASS \\
CxI & 0.8 & 0--10 & PASS \\
dSPN & 3.8 & 0--8 & PASS \\
iSPN & 4.2 & 0--8 & PASS \\
FSI & 7.8 & 2--30 & PASS \\
GPe & 60.7 & 20--80 & PASS \\
STN & 24.9 & 8--40 & PASS \\
GPi & 71.0 & 30--100 & PASS \\
Thalamus & 6.8 & 0--20 & PASS \\
\bottomrule
\end{tabular}
\end{table}

\subsection{Experiment 4: Lesion Studies}

\begin{table}[h]
\centering
\begin{tabular}{lcccl}
\toprule
Condition & GPi (Hz) & Th (Hz) & Clinical Match \\
\midrule
Healthy & 71.0 & 6.8 & Normal \\
Parkinson (D1 loss) & 81.0 & 3.4 & Bradykinesia (can't initiate) \\
Parkinson (D2 overactive) & 85.4 & 2.2 & Severe bradykinesia \\
Huntington (iSPN loss) & 50.4 & 15.0 & Hyperkinesia (involuntary moves) \\
STN lesion (DBS) & 0.0 & 46.2 & Therapeutic effect confirmed \\
\bottomrule
\end{tabular}
\end{table}

\end{document}
